\documentclass[11pt]{article}
\usepackage{amsmath,amssymb,amsthm,booktabs,hyperref}
\usepackage[margin=2.4cm]{geometry}
\usepackage{microtype}
\hypersetup{colorlinks=true,linkcolor=blue,citecolor=blue}
\newtheorem{theorem}{Theorem}
\newtheorem{lemma}{Lemma}
\newtheorem{proposition}{Proposition}
\newtheorem{corollary}{Corollary}
\theoremstyle{definition}
\newtheorem{definition}{Definition}
\newtheorem{hypothesis}{Hypothesis}
\newtheorem{conjecture}{Conjecture}
\newtheorem{numresult}{Numerical Result}
\theoremstyle{remark}
\newtheorem{remark}{Remark}
\title{Yield-Gap Reduction in Smooth Vector Liebig RFDEs:\\Finite-Time Control, Conditional Invariant Manifolds, and Local Spectral Persistence}
\author{Your Name}
\date{\today}

\begin{document}
\maketitle

\begin{abstract}
We study a smooth vector Liebig system whose binding and slack blocks form a weakly coupled retarded functional differential equation (RFDE). A uniform yield gap makes the off-limiting soft-minimum derivatives exponentially small, but this fact alone does not produce an invariant graph in the history phase space. We prove a mismatch-aware finite-time perturbation estimate and, under local exponential stability and stronger remainder control, a forward slack-tube estimate. We then state the precise compact normally attracting Banach-semiflow hypotheses under which a nearby invariant manifold would persist and admit a vertical representation over its perturbed projected base. Those geometric and spectral hypotheses are not verified for the concrete Liebig blocks, so the concrete graph remains conjectural. Independently of any graph, a simple transverse characteristic crossing persists with an $O(\varepsilon)$ shift. A nonlinear Hopf branch requires a separately matched RFDE Hopf theorem, adequate smoothness and nonlinear nondegeneracy; full orbital stability additionally requires stable complementary spectrum. At yield parity only the yield-gap certificate of small coupling is lost.
\end{abstract}

\section{Vector Liebig coupling and RFDE phase space}
\label{sec:object}

Let $i=1,\dots,n$ label companion blocks $X^i=(N_i,A_i,Z_i,E_i)\in\mathbb R^4$, with isolated-block vector fields inherited from~\cite{paperI}. The service entering the deficit is a smooth minimum
\[
S_{i,c}=S_{i,c}^{\max}\left[-\frac1\rho
\log\sum_m w_m\exp\left(-\rho\frac{\widetilde Y_{m,i,c}}{S_{i,c}^{\max}}\right)\right],
\]
with weights
\[
\pi_j=\frac{w_j e^{-\rho y_j}}{\sum_m w_m e^{-\rho y_m}}.
\]
Fix a constant delay $\tau>0$ and set
\[
X=C([ -\tau,0],\mathbb R^m),\qquad
Y=C([ -\tau,0],\mathbb R^n),\qquad
B=X\times Y,
\]
with the product supremum norm. After designating one block as binding, write the coupled RFDE as
\begin{equation}
\begin{aligned}
\dot x(t)&=F(x_t)+\varepsilon f(x_t,y_t),\\
\dot y(t)&=G(y_t)+\varepsilon g(x_t,y_t),
\end{aligned}
\label{eq:split}
\end{equation}
where $x_t(\theta)=x(t+\theta)$ and similarly for $y_t$. Throughout, all solution statements are localized to open domains and bounded trajectory enclosures on which the stated regularity and bounds hold.

\begin{definition}[Uniform yield gap]
\label{def:gap}
Component $k$ has a uniform yield gap $\Delta_y>0$ on a working region if
\[
\widetilde Y_k\le \widetilde Y_j-\Delta_y S^{\max}
\qquad (j\ne k)
\]
throughout that region.
\end{definition}

The soft-minimum formula gives
\[
\pi_j\le w_{\min}^{-1}e^{-\rho\Delta_y},\qquad j\ne k.
\]
Accordingly, the perturbation scale is written
\begin{equation}
\varepsilon=C_{\mathrm{gap}}e^{-\rho\Delta_y}+\varepsilon_{\mathrm{phys}}.
\label{eq:epsilon}
\end{equation}
The analysis below requires the numerical value of \eqref{eq:epsilon} to be small; a positive yield gap is one possible certificate, not the only one.

\section{Why the history cube is not an NHIM}
\label{sec:diagnosis}

The natural-looking product
\[
C([ -\tau,0],K_x)\times\{\widehat y_*\}
\label{eq:cube}
\]
is not justified as a compact smooth normally hyperbolic invariant manifold. If $K_x$ contains a nontrivial path (in particular, if it has nonempty interior), unrestricted continuous histories with values in $K_x$ need not be equicontinuous. For example, when $K_x=[-1,1]$, the histories $\phi_j(\theta)=\sin(j\theta)$ are bounded but have no uniformly convergent subsequence. Compactness of $K_x$ therefore does not imply compactness of its full history family. Degenerate targets, such as finite sets, are exceptions and do not supply the intended dynamical base.

Compactness would not by itself establish a manifold. A bounded absorbing set, attractor, or closed equicontinuous family need not be a $C^1$ embedded submanifold. Moreover, if a genuine binding manifold $\mathcal A_x\subset X$ exists, then
\[
T_{(\phi,\widehat y_*)}(\mathcal A_x\times\{\widehat y_*\})
=T_\phi\mathcal A_x\times\{0\},
\]
so every binding-history direction transverse to $T_\phi\mathcal A_x$, as well as every slack-history direction, belongs to the normal geometry. Stability of the isolated slack equilibrium controls only the latter part.

Finally, RFDE solution operators form a forward semiflow and need not admit ambient backward continuations. Persistence generally produces a nearby embedding
\[
\iota_\varepsilon(\phi)
=\bigl(\phi+u_\varepsilon(\phi),\widehat y_*+v_\varepsilon(\phi)\bigr),
\]
not automatically a vertical graph over the unchanged binding coordinate. These observations rule out a graph theorem based solely on a yield gap, slack stability, and bounded absorbing sets.

\section{Results that do not require an invariant graph}
\label{sec:proved}

\subsection{Finite-time perturbation control}

Write
\[
\mathcal H_0(x,y)=(F(x),G(y)),\qquad R(x,y)=(f(x,y),g(x,y)).
\]

\begin{theorem}[Finite-time tracking with initial mismatch]
\label{thm:track}
Let $z^\varepsilon$ solve \eqref{eq:split} and let $z^0$ solve the uncoupled product RFDE. Suppose both solutions exist on $[0,T]$ and remain in a common region on which
\[
\|\mathcal H_0(z_1)-\mathcal H_0(z_2)\|\le L\|z_1-z_2\|,
\qquad \|R(z)\|\le M.
\]
Then
\[
\sup_{0\le t\le T}\|z_t^\varepsilon-z_t^0\|_B
\le
\begin{cases}
 e^{LT}\|z_0^\varepsilon-z_0^0\|_B+
 \dfrac{|\varepsilon|M}{L}(e^{LT}-1),&L>0,\\[1.2ex]
 \|z_0^\varepsilon-z_0^0\|_B+|\varepsilon|MT,&L=0.
\end{cases}
\]
\end{theorem}

\begin{proof}
Set
\[
E(t)=\sup_{-\tau\le s\le t}|z^\varepsilon(s)-z^0(s)|.
\]
The two integral equations imply
\[
E(t)\le E(0)+L\int_0^tE(s)\,ds+|\varepsilon|Mt.
\]
Gronwall's inequality gives the stated bound. The estimate is finite-time and depends on the common enclosure; it does not imply all-time tracking.\qedhere
\end{proof}

\subsection{A forward slack tube}

\begin{hypothesis}[Local slack stability]
\label{hyp:slack}
The constant history $\widehat y_*$ satisfies $G(\widehat y_*)=0$. The slack linearization generates an evolution family with
\[
\|T_y(t)\|\le M_y e^{-\beta_y t},\qquad \beta_y>0.
\]
Near $\widehat y_*$, $G$ has a quadratic remainder (for example, $G$ is $C^{1,1}$ or $C^2$ there), and $g$ is bounded on the working region.
\end{hypothesis}

\begin{theorem}[Confined forward slack tube]
\label{thm:tube}
Under Hypothesis~\ref{hyp:slack}, fix any $\beta'\in(0,\beta_y)$. There are $\delta>0$, $\varepsilon_0>0$, and $C\ge1$ such that, if
\[
\|y_0-\widehat y_*\|_Y\le\delta,
\]
and the binding history remains in the declared bounded working region for all times considered, then for $|\varepsilon|\le\varepsilon_0$,
\[
\|y_t-\widehat y_*\|_Y
\le C\left(e^{-\beta't}\|y_0-\widehat y_*\|_Y+|\varepsilon|\right).
\]
\end{theorem}

\begin{proof}
Use the RFDE variation-of-constants formula in its fundamental-solution (or equivalent sun--star) formulation. Exponential convolution bounds the forcing by $C|\varepsilon|$, while the nonlinear remainder is bounded by $C\|y_t-\widehat y_*\|_Y^2$. A bootstrap on the ball of radius $\delta$ absorbs the quadratic term and yields the estimate with any fixed $\beta'<\beta_y$. The binding-confinement assumption is essential.\qedhere
\end{proof}

An initially $O(\varepsilon)$ slack error needs no logarithmic entrance time. An order-one local error reaches an $O(\varepsilon)$ tube only after a transient of order $|\log\varepsilon|$.

\subsection{Compact history subsets and equilibria}

\begin{lemma}[A compact history class]
\label{lem:equi}
A closed family of histories that is uniformly bounded and equi-Lipschitz is compact in the supremum norm. It is positively invariant only if the RFDE vector field and boundary barriers preserve the stated value and Lipschitz bounds.
\end{lemma}

\begin{proof}
Compactness is Arzel\`a--Ascoli; positive invariance is a separate dynamical condition.\qedhere
\end{proof}

This lemma supplies useful compact subsets but not a manifold or a graph domain.

\begin{lemma}[Equilibrium continuation]
\label{lem:eq}
Suppose $F(\widehat x_*)=0$, $G(\widehat y_*)=0$, and the equilibrium Jacobian is invertible. Equivalently, under the convention
\[
\Delta(\lambda)=\lambda I-L(e^{\lambda\cdot}I),
\]
assume $\Delta_x(0)$ and $\Delta_y(0)$ are invertible. Then a unique local equilibrium branch $(x_*(\varepsilon),y_*(\varepsilon))$ exists and
\[
|x_*(\varepsilon)-x_*|+|y_*(\varepsilon)-y_*|=O(|\varepsilon|).
\]
\end{lemma}

\begin{proof}
Constant equilibria are zeros of the finite-dimensional map
\[
\mathcal F(x,y,\varepsilon)=
\begin{pmatrix}
F(\widehat x)+\varepsilon f(\widehat x,\widehat y)\\
G(\widehat y)+\varepsilon g(\widehat x,\widehat y)
\end{pmatrix}.
\]
Its state derivative at $\varepsilon=0$ has diagonal blocks $-\Delta_x(0)$ and $-\Delta_y(0)$. The finite-dimensional implicit-function theorem gives the branch and estimate.\qedhere
\end{proof}

A hyperbolic periodic orbit is another possible compact finite-dimensional object, but its persistence requires an exact RFDE periodic-orbit theorem or a fully constructed Poincar\'e return map. Neither an equilibrium nor a periodic orbit gives a graph over a region of histories.

\section{Conditional compact-NAIM persistence}
\label{sec:conditional}

The correct abstract route is a compact normally attracting invariant-manifold theorem for Banach-space semiflows, such as an applicable result in Bates, Lu, and Zeng~\cite{blz1998}. It is not finite-dimensional Fenichel theory and it is not an inversion of the ambient RFDE time map.

\begin{hypothesis}[Compact Banach-semiflow NAIM data]
\label{hyp:naim}
Assume the following data are verified for the uncoupled RFDE.
\begin{enumerate}
\item There is a compact finite-dimensional $C^1$ embedded invariant manifold $\mathcal A_x\subset X$, without boundary (or with the precise inflowing/overflowing boundary required by the selected theorem), and the restricted dynamics have the flow property required by that theorem.
\item The product $\mathcal M_0=\mathcal A_x\times\{\widehat y_*\}$ has a uniformly complemented invariant splitting
\[
TB|_{\mathcal M_0}=E^s\oplus T\mathcal M_0
\]
(or the theorem's full $E^s\oplus T\mathcal M_0\oplus E^u$ splitting), containing every binding-transverse and slack direction.
\item The derivative cocycle obeys the selected theorem's normal and tangent estimates. For a time-$T_0$ map, a representative $C^1$ bunching condition is
\[
\sup_{p\in\mathcal M_0}
\|DP_0(p)|_{E_p^s}\|
\left\|(DP_0(p)|_{T_p\mathcal M_0})^{-1}\right\|<1.
\]
Higher regularity requires the exact higher-order conorm inequalities, not merely signs of spectral abscissae.
\item A uniform tubular neighborhood, projections, localization, semiflow regularity, and boundary conditions match the selected theorem.
\item On the required neighborhood and compact positive-time interval, the perturbed semiflow maps are $C^1$-$O(\varepsilon)$ close to the unperturbed maps, and the selected theorem (or its graph-transform proof) supplies Lipschitz dependence of the persistent embedding on that perturbation. Quantitative higher-order closeness is assumed only when separately proved.
\end{enumerate}
\end{hypothesis}

\begin{proposition}[Conditional persistence and vertical representation]
\label{prop:conditional}
If Hypothesis~\ref{hyp:naim} matches all hypotheses of the selected Banach-semiflow persistence theorem, then that theorem supplies a nearby locally positively invariant $C^1$ manifold
\[
\mathcal M_\varepsilon=\iota_\varepsilon(\mathcal A_x),
\qquad
\|\iota_\varepsilon-\iota_0\|_{C^1}=O(|\varepsilon|),
\]
with the theorem's local normal-attraction conclusion. If, additionally,
\[
q_\varepsilon=\pi_x\circ\iota_\varepsilon
\]
is a $C^1$ embedding, then $q_\varepsilon$ is a diffeomorphism onto
\[
\mathcal A_{x,\varepsilon}=q_\varepsilon(\mathcal A_x),
\]
and
\[
h_\varepsilon=
\pi_y\circ\iota_\varepsilon\circ q_\varepsilon^{-1}
\]
gives
\[
\mathcal M_\varepsilon
=\{(\phi,h_\varepsilon(\phi)): \phi\in\mathcal A_{x,\varepsilon}\}.
\]
\end{proposition}

\begin{proof}
The persistence and attraction conclusions are precisely those of the matched external theorem. The projection statement is separate. Since $\mathcal A_x$ is compact, a sufficiently small $C^1$ perturbation of its inclusion remains an embedding: local injectivity follows from the tangent conorm, while global injectivity follows by separating near pairs in finitely many charts and far pairs by compactness. The formula for $h_\varepsilon$ then follows by reparametrization.\qedhere
\end{proof}

The graph is over the perturbed base $\mathcal A_{x,\varepsilon}$, not over the old coordinate copy of $\mathcal A_x$ and never over the full history cube.

\begin{conjecture}[Concrete A021 compact-NAIM application]
\label{conj:graph}
For a concrete A021 binding block, there exists a compact finite-dimensional binding manifold $\mathcal A_x$ and a complete invariant splitting satisfying Hypothesis~\ref{hyp:naim}, including quantitative inverse-tangent domination and the projection-embedding condition.
\end{conjecture}

The companion working core supplies a concrete candidate and substantial numerical evidence, but not yet the complete theorem package; the status is made precise next.

\subsection{Companion-core crosswalk and numerical periodic candidate}
\label{sec:c4-crosswalk}

The turnover-corrected gated four-state working core of~\cite{paperI}, with state $(N,A,Z,E)$, is the natural architectural candidate for the abstract binding functional $F^k$ in~\eqref{eq:split}. Under the companion's dynamic-target, large-reservoir closure, its five-state extension has a driven detritus coordinate that does not feed back into $(N,A,Z,E)$; the fixed-target closure does not share that projection. This identification does not specify the slack blocks or the coordinate-level multi-block residuals $f$ and $g$.

\begin{numresult}[Gated C4 periodic candidate at $\tau_x=4.5$]
\label{num:c4-cycle}
For the Candidate-A gated four-state working core at delay $\tau_x=4.5$~yr, a method-of-steps RK4 reproduction converges to a large-cycle candidate in the lower bistable window with period
\[
P_x\approx370.95~\mathrm{yr}.
\]
Over the final computed cycles,
\[
N\in[45.69208,94.87305],\qquad
A\in[834.58311,943.05308],
\]
\[
Z\in[0.001534,0.678093],\qquad
E\in[0.354363,20.082770].
\]
The full argument of the outer memory floor remains positive, with minimum approximately $1.4755\times10^{-3}$; hence this orbit lies in a locally smooth region of the vector field.

The complete finite-history monodromy matrices at three grid levels give
\begin{center}
\begin{tabular}{@{}rrrr@{}}
\toprule
$\Delta t$ (yr) & history dimension & phase multiplier & leading nontrivial multiplier\\
\midrule
$0.25$ & $76$  & $0.98687854$ & $0.68774849$\\
$0.10$ & $184$ & $0.99774865$ & $0.68770289$\\
$0.05$ & $364$ & $1.00136091$ & $0.68768669$\\
\bottomrule
\end{tabular}
\end{center}
A linear step-size extrapolation gives $\mu_s\approx0.68767164$. Twice the finest two-level difference gives the conservative empirical discretization envelope
\[
\mu_s\in[0.68763924,0.68770405],
\]
corresponding to the leading numerical normal exponent
\[
\beta_x\approx1.0094\times10^{-3}~\mathrm{yr}^{-1}.
\]
This interval is not an outward-rounded enclosure of the continuum monodromy operator. It supplies strong numerical evidence, not a proof of algebraic simplicity of the phase multiplier or of the complete continuum Floquet spectrum.
\end{numresult}

\paragraph{Fourier seed and computer-assisted-proof specification.}
A phase-corrected Fourier seed with 512 positive and 512 negative modes per state has period $370.9311778463889$~yr, endpoint mismatch $4.90\times10^{-7}$, and RFDE spectral residual at most $6.80\times10^{-6}$. The statewise maximum residuals are approximately
\[
(8.26\times10^{-10},\ 6.79\times10^{-6},\ 1.05\times10^{-8},\ 1.48\times10^{-6}).
\]
The full memory-floor argument has numerical minimum $1.4755423\times10^{-3}$ on this seed. A five-module computer-assisted specification is formulated: a weighted-Fourier radii polynomial for the orbit, a continuum Floquet enclosure with Riesz projections, a validated slack root/semigroup calculation, a continuum product-bunching test with target $q_{40}<1/4$, and the remaining coupling/semiflow/projection gates. Outward-rounded execution is pending: no Fourier weight, tail bound, radii-polynomial constants, continuum multiplier disks, or operator-error constants are asserted. The seed and specification therefore do not alter the conditional status of Proposition~\ref{prop:conditional} or Conjecture~\ref{conj:graph}.

The phase-tangent history norm ratio on the reproduced orbit is approximately $M_c=4.5536$. A finite-discrete phase-projection norm near $547$ indicates strong nonnormal transients; spectral radius alone therefore cannot verify bunching.

For a minimal numerical product scaffold, take an identical Candidate-A C4 slack block at its equilibrium with delay $\tau_y=10$~yr, on the common maximal history interval $[-10,0]$. Determinant refinement gives the rightmost pair
\[
\lambda_{y,\pm}
=-0.00052673009564114
\pm0.0220846350193287i.
\]
An analytic exterior exclusion combined with high-precision argument-principle counts isolates this as the unique rightmost pair at numerical-certification status. The provisional product rate is therefore
\[
\beta=\min(\beta_x,\beta_y)
\approx5.2673\times10^{-4}~\mathrm{yr}^{-1},
\]
controlled by the slack equilibrium.

For a time-$T$ map, a representative $C^1$ target is
\[
M_sM_c e^{-\beta T}<1.
\]
Direct induced-supremum-norm calculations on the finite history discretizations include the large transient prefactors rather than replacing them by one. Extrapolating the slack method-of-lines norms in inverse history resolution gives the numerical products
\[
q_{30}\approx0.952,
\qquad q_{35}\approx0.338,
\qquad q_{40}\approx0.116
\]
at $30$, $35$, and $40$ binding periods, respectively. Thus discrete product bunching closes only marginally at thirty periods and robustly by thirty-five periods. These are finite-discretization operator-norm estimates, not a continuum domination proof.

\begin{remark}[Status of the numerical crosswalk]
The selected binding cycle and slack equilibrium define a concrete uncoupled numerical scaffold
\[
\Gamma_x\times\{\widehat y_*\}.
\]
A theorem still requires a validated continuum periodic orbit and Floquet enclosure, invariant history-space projections and prefactors, source-derived slack/coupling functionals $G,f,g$, a uniform $C^1$ perturbation tube, and an exact Banach-semiflow theorem match. Any persistent vertical representation would be over a perturbed projected closed curve, not over a region of histories. Conjecture~\ref{conj:graph} therefore remains open and Proposition~\ref{prop:conditional} remains conditional.
\end{remark}

\section{Direct characteristic crossing and conditional nonlinear Hopf}
\label{sec:hopf}

The spectral route below is independent of an invariant graph. Let $\mu$ be a physical bifurcation parameter. If $\mu$ changes the delay, the RFDE family must first be formulated smoothly on a fixed phase space, or an exact varying-delay theorem must be used.

\begin{theorem}[Persistence of a simple characteristic crossing]
\label{thm:crossing}
Assume a $C^1$ equilibrium branch $z_*(\mu,\varepsilon)$ and a characteristic function
\[
D(\lambda,\mu,\varepsilon)
\]
that is holomorphic in $\lambda$ and $C^1$ in $(\mu,\varepsilon)$. At $(\mu_*,0)$ assume:
\begin{enumerate}
\item $D$ has a simple pair $\pm i\omega_*$, $\omega_*>0$;
\item no other characteristic root lies on the imaginary axis, and the full complementary spectrum is locally separated from the pair;
\item the crossing is transverse:
\[
\partial_\mu\operatorname{Re}\lambda(\mu_*,0)\ne0.
\]
\end{enumerate}
Then there are $C^1$ functions
\[
\mu_*(\varepsilon)=\mu_*+O(\varepsilon),
\qquad
\omega(\varepsilon)=\omega_*+O(\varepsilon)
\]
such that
\[
D(i\omega(\varepsilon),\mu_*(\varepsilon),\varepsilon)=0.
\]
\end{theorem}

\begin{proof}
Apply the real implicit-function theorem to
\[
(\omega,\mu,\varepsilon)\longmapsto
\begin{pmatrix}
\operatorname{Re}D(i\omega,
\mu,
\varepsilon)\\
\operatorname{Im}D(i\omega,
\mu,
\varepsilon)
\end{pmatrix}.
\]
Simplicity and transversality make the Jacobian in $(\omega,\mu)$ invertible. Standard RFDE root localization and spectral perturbation preserve isolation of the pair.\qedhere
\end{proof}

Theorem~\ref{thm:crossing} is a spectral theorem, not a periodic-orbit theorem.

\begin{proposition}[Conditional nonlinear Hopf conclusion]
\label{prop:hopf}
Suppose, in addition to Theorem~\ref{thm:crossing}, that the RFDE family is represented on a fixed phase space and satisfies every hypothesis of a matched RFDE Hopf theorem, including its smoothness, nonresonance, equilibrium, center-dimension, and transversality assumptions. Then that theorem yields its stated local periodic branch, unique only in the theorem's local sense and up to phase.

If a first Lyapunov coefficient is used to classify the branch, assume the regularity and nonlinear nondegeneracy required by the matched normal-form theorem and verify
\[
\ell_1(0)\ne0.
\]
Continuity of the correctly normalized coefficient in the theorem's jet and spectral-projection topology then preserves its sign for small $\varepsilon$. Full orbital asymptotic stability in $B$ additionally requires every complementary characteristic root of the full coupled RFDE to lie strictly in the left half-plane. No local Hopf theorem implies persistence of a global fold, isola, or connecting orbit.
\end{proposition}

\section{Yield parity}
\label{sec:parity}

\begin{proposition}[Loss of the yield-gap certificate]
\label{prop:parity}
At yield parity, $\Delta_y=0$, equation~\eqref{eq:epsilon} becomes
\[
\varepsilon=C_{\mathrm{gap}}+\varepsilon_{\mathrm{phys}}.
\]
Thus the yield gap no longer makes the coupling exponentially small. Every conclusion requiring small $\varepsilon$ loses this particular certificate.
\end{proposition}

Parity does not prove that $\varepsilon$ is numerically large: $C_{\mathrm{gap}}$ or $\varepsilon_{\mathrm{phys}}$ may be independently small. Nor does parity exclude a reduction based on another spectral gap, symmetry, averaging, or timescale separation.

\section{Concrete verification obligations}
\label{sec:checklist}

Promotion of Conjecture~\ref{conj:graph} or Proposition~\ref{prop:hopf} for an A021 block requires:
\begin{enumerate}
\item an explicit compact finite-dimensional invariant binding manifold, with dimension, charts, and boundary status;
\item the restricted base-flow property required by the selected theorem;
\item a complete invariant splitting containing every binding-transverse and slack mode;
\item quantitative tangent norm and inverse-tangent bounds, and every stable or unstable normal rate;
\item the exact domination inequalities of the selected theorem;
\item a localized perturbation bound in the required $C^1$ (or stronger) topology and a common trajectory enclosure;
\item tubular geometry and, for a vertical graph, global projected-embedding verification;
\item an independently verified uniform yield gap if it is used as the smallness certificate;
\item for Hopf, a fixed-phase-space parameter family, equilibrium continuation, the complete characteristic spectrum, simplicity, transversality, and theorem-matched smoothness;
\item for branch criticality and stability, a correctly normalized nonzero Lyapunov coefficient and strict stability of the full complementary spectrum.
\end{enumerate}

\section{Conclusion}

The yield gap supplies a quantitatively useful perturbation scale, but it does not manufacture compactness, manifold structure, or normal hyperbolicity in RFDE history space. The unconditional reduction result is finite-time perturbation control. A forward slack tube follows under stronger local stability and confinement hypotheses. A compact normally attracting invariant manifold and its vertical representation are available only conditionally through an exactly matched Banach-semiflow theorem, and the concrete A021 hypotheses remain open. Local characteristic crossing persists directly, without a graph; nonlinear Hopf and orbital stability require separate theorem matching and additional nonlinear and spectral data. At parity only the yield-gap smallness mechanism is lost.

\begin{thebibliography}{9}
\bibitem{paperI}
Y.\ Name, ``Scarcity-driven capital liquidation\ldots,'' companion manuscript, 2026.

\bibitem{hale1993}
J.\ K.\ Hale and S.\ M.\ Verduyn Lunel,
\emph{Introduction to Functional Differential Equations},
Springer, 1993.

\bibitem{blz1998}
P.\ W.\ Bates, K.\ Lu, and C.\ Zeng,
\emph{Existence and Persistence of Invariant Manifolds for Semiflows in Banach Space},
Memoirs of the American Mathematical Society, vol.~135, no.~645, 1998.
\end{thebibliography}

\end{document}
